% ============================================
% CAPÍTULO 4 — Obra Completa (Referência)
% ============================================

\chapter*{A Obra Completa — 70 Fragmentos}
\addcontentsline{toc}{chapter}{A Obra Completa — 70 Fragmentos}

\epigraph{Este arquivo contém aquilo que o Hospital Colônia não conseguiu queimar.}{--- O Arquivista}

\noindent A versão impressa que você tem em mãos contém uma seleção editorial dos fragmentos mais representativos da obra. A versão \textbf{completa}, com todos os 70 fragmentos organizados em ordem linear, está disponível no formato \texttt{Markdown} no arquivo:

\begin{center}
\vspace{0.5cm}
\fbox{\parbox{0.7\textwidth}{\centering\texttt{projetos/molambudos/molambudos.md}}}
\vspace{0.5cm}
\end{center}

\section*{Estrutura Completa da Obra}

A obra completa contém \textbf{70 fragmentos} distribuídos em 4 categorias:

\begin{itemize}
    \item \textbf{26 MEM} — Memórias do Paciente 1.260 (Joaquim Antônio Correia)
    \item \textbf{19 DOC} — Documentos oficiais, laudos e registros históricos
    \item \textbf{13 CONT} — A contaminação progressiva do leitor
    \item \textbf{12 LUC} — A investigação forense da Dra. Lúcia Menezes
\end{itemize}

\vspace{0.3cm}

\begin{tabularx}{\textwidth}{|X|X|X|}
\hline
\textbf{Período} & \textbf{Fragmentos} & \textbf{Conteúdo} \\
\hline
1915--1928 & MEM-01 a MEM-07, DOC-17, DOC-18 & O sertão, a seca, a vala, os currais do governo, a mãe levada \\
\hline
1928--1935 & MEM-08 a MEM-16, DOC-02, DOC-10 & A chegada ao Colônia, os primeiros procedimentos, as vozes \\
\hline
1935--1979 & MEM-17 a MEM-26, DOC-01, DOC-03 a DOC-09, DOC-11 a DOC-16 & O diário de Joaquim, os laudos, o diário do Dr. Oliveira \\
\hline
1979--2026 & DOC-19, DOC-08, LUC-01 a LUC-12 & Exumação, investigação, a contaminação de Lúcia \\
\hline
\text{Ciclo} & CONT-01 a CONT-13 & O olho amarelo, o Rasga Mortalha, o paciente 1.263 \\
\hline
\end{tabularx}

\vspace{0.5cm}

\section*{Escala de Contaminação do Leitor}

\begin{longtable}{|c|p{10cm}|}
\hline
\textbf{Pontos} & \textbf{Diagnóstico} \\
\hline
0--2 & Leitor resiliente. A contaminação não pegou. Ou você está mentindo. \\
\hline
3--5 & Leitor exposto. Sintomas iniciais. Recomenda-se pausa na leitura. \\
\hline
6--8 & Leitor contaminado. O ciclo está em andamento. O olho amarelo aparecerá em breve. \\
\hline
9--10 & Paciente 1.263. Bem-vindo. A criatura agradece sua atenção. \\
\hline
\end{longtable}

\vspace{0.5cm}

\section*{Estatísticas da Obra Completa}

\begin{itemize}
    \item \textbf{Total de fragmentos:} 70
    \item \textbf{Total de linhas:} 7.511
    \item \textbf{Total de caracteres:} ~225.000
    \item \textbf{Média por fragmento:} ~107 linhas
    \item \textbf{Período coberto:} 1915--2026 (111 anos)
    \item \textbf{Personagens documentados:} 40+
    \item \textbf{Procedimentos psiquiátricos descritos:} 6 tipos (ECT, lobotomia, insulinoterapia, contenção, banhos térmicos, camisa-de-força)
    \item \textbf{Marcas de violência no corpo:} 59 anos acumulados
\end{itemize}

\vspace{0.5cm}

\section*{Atos de Horror Progressivo}

\begin{enumerate}
    \item \textbf{Horror Visceral} (MEM-01 a MEM-07) — Seca, fome, morte do gado, vala de retirantes. A prosa poética e brutal de um menino de 8 anos.
    \item \textbf{Horror Institucional} (MEM-08 a MEM-26, DOC-01 a DOC-16) — O sistema psiquiátrico como máquina de moer corpos. Eletrochoque, lobotomia, insulinoterapia. A voz clínica que se desintegra.
    \item \textbf{Horror Metafísico} (LUC-01 a LUC-12, CONT-01 a CONT-13) — A contaminação que salta da página. O olho amarelo. O Rasga Mortalha. O ciclo de pacientes numerados.
\end{enumerate}

\vspace{0.8cm}

\noindent\rule{\textwidth}{0.5pt}

\vspace{0.3cm}

\begin{center}
\textit{O ciclo recomeça a cada leitor.} \\
\textit{Ao final, você é o Paciente 1.263.} \\
\textit{Bem-vindo ao arquivo.} \\
\vspace{0.5cm}
\rule{0.3\textwidth}{0.5pt} \\
\vspace{0.3cm}
\textit{Se você está lendo isto, você já é o 1.263.} \\
\textit{O cheiro já começou?} \\
\textit{--- Lúcia Menezes, Paciente 1.262}
\end{center}

\endinput
